\documentclass[../../main.tex]{subfiles}
\begin{document}

{\bf [解]}

% TODO: S,Tの概要
\begin{figure}[htb]
  \centering
  % 俯角 65°，方位角 120°
  \tdplotsetmaincoords{65}{120}
  \begin{tikzpicture}[tdplot_main_coords, scale=2.2, line join=round, >=stealth, every node/.style={font=\small}]

    % --- 1. 定数・座標設定 ---
    % 点 P = (0.9, 0.9, 1.4) (OP^2 = 0.81 + 0.81 + 1.96 = 3.58 > 1)
    \coordinate (O) at (0, 0, 0);
    \coordinate (A) at (0, 0, -1.0);
    \coordinate (P) at (0.9, 0.9, 1.4);

    % OP の中点 M (球面 T の中心)
    \coordinate (M) at (0.45, 0.45, 0.7);

    % 平面 L: Xx + Yy + Zz = 1
    % 原点から平面 L への垂線の足 H_L: H_L = (1 / |OP|^2) * P
    \pgfmathsetmacro{\opSq}{0.9*0.9 + 0.9*0.9 + 1.4*1.4} % 3.58
    \pgfmathsetmacro{\scaleH}{1.0 / \opSq}               % ≈ 0.2793
    \pgfmathsetmacro{\xHL}{\scaleH * 0.9}                % ≈ 0.2514
    \pgfmathsetmacro{\yHL}{\scaleH * 0.9}                % ≈ 0.2514
    \pgfmathsetmacro{\zHL}{\scaleH * 1.4}                % ≈ 0.3911
    \coordinate (HL) at (\xHL, \yHL, \zHL);

    % 点 P から平面 L への垂線の足 Q (= HL)
    \coordinate (Q) at (HL);

    % 点 A(0, 0, -1) から平面 L への垂線の足 R:
    \pgfmathsetmacro{\scaleR}{(1.4 + 1.0) / \opSq}       % 2.4 / 3.58 ≈ 0.6704
    \coordinate (R_pt) at ({\scaleR*0.9}, {\scaleR*0.9}, {-1.0 + \scaleR*1.4});

    % --- 2. 座標軸 ---
    \draw[dashed, gray!40] (-1.3, 0, 0) -- (0, 0, 0);
    \draw[dashed, gray!40] (0, -1.3, 0) -- (0, 0, 0);
    \draw[dashed, gray!40] (0, 0, -1.4) -- (0, 0, 0);

    \draw[->, thick, gray!80!black] (0, 0, 0) -- (1.8, 0, 0) node[below left=1pt] {$x$};
    \draw[->, thick, gray!80!black] (0, 0, 0) -- (0, 1.8, 0) node[right=2pt] {$y$};
    \draw[->, thick, gray!80!black] (0, 0, 0) -- (0, 0, 2.0) node[above=2pt] {$z$};
    \node[below left=2pt] at (O) {$O$};

    % --- 3. 球面 S (単位球) の大円 ---
    \draw[thick, gray!55] plot[domain=0:360, samples=60] ({cos(\x)}, {sin(\x)}, 0);
    \draw[dashed, gray!45] plot[domain=0:360, samples=60] ({cos(\x)}, 0, {sin(\x)});
    \draw[dashed, gray!45] plot[domain=0:360, samples=60] (0, {cos(\x)}, {sin(\x)});
    \node[gray!70!black, font=\scriptsize, below left] at (-0.6, -0.6, 0) {球面 $S$};

    % --- 4. 球面 T (OP を直径とする球面) の円周 ---
    \pgfmathsetmacro{\radT}{sqrt(\opSq)/2.0}
    \draw[dashed, teal!65, thick]
    plot[domain=0:360, samples=50] ({0.45 + \radT*cos(\x)}, {0.45 + \radT*sin(\x)}, 0.7);
    \node[teal!90!black, font=\scriptsize] at ($(M) + (0.5, 0.5, 0.7)$) {球面 $T$};

    % --- 5. 平面 L (交円を含む平面パッチ) ---
    \coordinate (P_L1) at ($(HL) + (-0.9, 0.9, 0)$);
    \coordinate (P_L2) at ($(HL) + (0.6, 0.9, -0.77)$);
    \coordinate (P_L3) at ($(HL) + (0.9, -0.9, 0)$);
    \coordinate (P_L4) at ($(HL) + (-0.6, -0.9, 0.77)$);

    \fill[blue!15, opacity=0.55] (P_L1) -- (P_L2) -- (P_L3) -- (P_L4) -- cycle;
    \draw[line width=1.2pt, blue!75!black] (P_L1) -- (P_L2) -- (P_L3) -- (P_L4) -- cycle;
    \node[blue!90!black, font=\footnotesize\bfseries, above right] at (P_L2) {平面 $L$};

    % S と T の交円（成分ごとに直接計算して安全にプロット）
    % 半径 r_c = 0.849, 基底ベクトル u = (0.7071, -0.7071, 0), v = (-0.516, -0.516, 0.68)
    \draw[very thick, blue!85!black, domain=0:360, samples=70]
    plot ({ \xHL + 0.849*cos(\x)*0.7071 - 0.849*sin(\x)*0.516 },
      { \yHL - 0.849*cos(\x)*0.7071 - 0.849*sin(\x)*0.516 },
    { \zHL + 0.849*sin(\x)*0.68 });

    % --- 6. 点 P, Q, A, R の関係線 ---
    % 直線 OP (直径、平面 L に直交)
    \draw[thick, gray!70] (O) -- (Q);
    \draw[line width=1.8pt, red!85!black] (Q) -- (P);
    \node[red!90!black, font=\scriptsize\bfseries, right] at ($ (Q)!0.5!(P) $) {$|PQ|$};

    % 点 A(0,0,-1) と平面 L への垂線 AR
    \draw[line width=1.8pt, purple!85!black] (A) -- (R_pt);
    \node[purple!90!black, font=\scriptsize\bfseries, right] at ($ (A)!0.5!(R_pt) $) {$|AR|$};

    % 垂線の直角マーク
    \draw[thin, gray!70]
    ($ (Q) + (0.08, -0.08, 0) $) --
    ($ (Q) + (0.08, -0.08, 0) + (0.06, 0.06, 0.09) $) --
    ($ (Q) + (0.06, 0.06, 0.09) $);

    % --- 7. 主要点のプロットとラベル ---
    \fill[black]         (O)    circle (1.0pt);
    \fill[black]         (A)    circle (1.2pt) node[below left=1pt] {$A(0,0,-1)$};
    \fill[red!85!black]  (P)    circle (1.4pt) node[above right=1pt, red!95!black] {$P(X,Y,Z)$};
    \fill[teal!85!black] (M)    circle (0.9pt) node[above left=1pt, font=\scriptsize] {$M$};
    \fill[blue!90!black] (Q)    circle (1.2pt) node[above left=1pt, font=\scriptsize] {$Q$};
    \fill[purple!90!black] (R_pt) circle (1.2pt) node[above right=1pt, font=\scriptsize] {$R$};

  \end{tikzpicture}
  \caption{球面 $S, T$，点 $P, A$ および平面 $L$ への垂線 $PQ, AR$ の空間配置}
  \label{fig:spheres_plane_orthogonal_projection_3d}
\end{figure}

$A(0,0,-1)$だから$|OA|=1$であり，よって球面$S$は半径$1$でその方程式は
\begin{align}
  S: f(x,y,z) = x^2+y^2+z^2 = 1 \label{eq:1}
\end{align}
である．
$OP$を直径とする球面を$T$とすると，$T$の中心は$OP$の中点$M$で，$P(X,Y,Z)$とおくと$M(X/2,Y/2,Z/2)$であるから$T$の方程式は
\begin{align}
  & T:\left(x-\dfrac{X}{2}\right)^2 + \left(y-\dfrac{Y}{2}\right)^2
  +\left(z-\dfrac{Z}{2}\right)^2
  = \left(\dfrac{X}{2}\right)^2 + \left(\dfrac{Y}{2}\right)^2 + \left(\dfrac{Z}{2}\right)^2 \\
  \therefore
  & T: g(x,y,z) = x^2-Xx+y^2-Yy+z^2-Zz=0 \label{eq:2}
\end{align}
である．
点$P$は$S$の外部だから
\begin{align}
  f(X,Y,Z)>0 \label{eq:3}
\end{align}
が成立する．
このもとで$T,S$は必ず交円を持つ．ここで$k\in\mathbb R$に対し
\begin{align}
  f(x,y,z)+kg(x,y,z)=0 \label{eq:4}
\end{align}
で表される方程式は$f=g=0$を解に持つから$S,T$の交円を通る．
$k=-1$の時\cref{eq:4}は平面の方程式になり，またLはただ1通りに定まるからこれが$L$に他ならない．したがって，\cref{eq:1,eq:2}を代入して
\begin{align}
  L:Xx+Yy+Zz-1=0 \label{eq:5}
\end{align}
である．$|PQ|$と$|AR|$はそれぞれ点$P, A$と平面$L$との距離だから，点と平面の距離の公式から
\begin{align}
  \begin{cases}
    |PQ|&=\dfrac{|X^2+Y^2+Z^2-1|}{\sqrt{X^2+Y^2+Z^2}}
    =\dfrac{X^2+Y^2+Z^2-1}{\sqrt{X^2+Y^2+Z^2}}\quad(\because\text{\cref{eq:3}})\\
    |AR|&=\dfrac{|{-Z-1}|}{\sqrt{X^2+Y^2+Z^2}}=\dfrac{|Z+1|}{\sqrt{X^2+Y^2+Z^2}}
  \end{cases}
\end{align}
である．

だから題意の条件$|PQ|\le|AR|$に代入して，
\begin{align}
  X^2+Y^2+Z^2-1\le|Z+1|
\end{align}
である．絶対値を外して
\begin{align}
  &
  \begin{cases}
    X^2+Y^2+Z^2-1 \le Z+1  & (Z+1 \ge 0) \\
    X^2+Y^2+Z^2-1 \le -Z-1 & (Z+1 \le 0)
  \end{cases} \\
  \therefore
  &
  \begin{cases}
    X^2+Y^2+Z^2-Z \le 2  & (Z \ge -1) \\
    X^2+Y^2+Z^2+Z \le 0  & (Z \le -1)
  \end{cases} \\
  \therefore
  &
  \begin{cases}
    X^2+Y^2+\left(Z-\dfrac{1}{2}\right)^2  \le \left(\dfrac{3}{2}\right)^2 & (Z \ge -1) \\
    X^2+Y^2+\left(Z+\dfrac{1}{2}\right)^2  \le \left(\dfrac{1}{2}\right)^2 & (Z \le -1)
  \end{cases}
  \therefore
  &
  \begin{cases}
    X^2+Y^2+\left(Z-\dfrac{1}{2}\right)^2  \le \left(\dfrac{3}{2}\right)^2 & (Z \ge -1) \\
    (X,Y,Z)=(0,0,-1)
  \end{cases}
\end{align}
である．後者は前者に含まれるから
\begin{align}
  X^2+Y^2+\left(Z-\dfrac{1}{2}\right)^2  \le \left(\dfrac{3}{2}\right)^2 & (Z \ge -1) \label{eq:6}
\end{align}
である．

従って\cref{eq:3,eq:6}が求める条件で
\begin{align}
  \begin{cases}
    X^2+Y^2+Z^2 > 1 \\
    X^2+Y^2+\left(Z-\dfrac{1}{2}\right)^2 \le \left(\dfrac{3}{2}\right)^2
  \end{cases} \label{eq:7}
\end{align}
である．これを図示すると下図の斜線部のようになる（$XZ$平面の切断面）．

% TODO: XZでの\cref{eq:7}の切断面
\begin{figure}[htb]
  \centering
  \begin{tikzpicture}[scale=2.2, >=stealth, every node/.style={font=\small}]

    % --- 1. 定数・座標パラメータ ---
    \coordinate (O) at (0, 0);          % 原点 O(0, 0)
    \coordinate (M) at (0, 0.5);        % 外側球の中心 M(0, 1/2)
    \coordinate (A) at (0, -1.0);       % 接点 A(0, -1)
    \coordinate (Top) at (0, 2.0);      % 外側球の最上点 (0, 2)

    % --- 2. 求める領域のハッチング塗りつぶし (even odd rule) ---
    % 外円 (半径 1.5, 中心 (0, 0.5)) から 内円 (半径 1.0, 中心 (0, 0)) を除外
    \fill[pattern=north east lines, pattern color=blue!65, opacity=0.75, even odd rule]
    (M) circle (1.5) (O) circle (1.0);

    % --- 3. 境界円の描画 ---
    % 外側の円 (境界を含むため実線)
    \draw[very thick, blue!85!black] (M) circle (1.5);
    \node[blue!90!black, font=\footnotesize\bfseries, above right] at (1.1, 1.4)
    {$X^2 + \left(Z - \frac{1}{2}\right)^2 = \left(\frac{3}{2}\right)^2$};

    % 内側の円 (単位円 S の大円, 境界を含まないため破線)
    \draw[thick, dashed, red!80!black] (O) circle (1.0);
    \node[red!90!black, font=\footnotesize\bfseries, above left] at (-0.6, 0.6)
    {$S: X^2 + Z^2 = 1$};

    % --- 4. 座標軸 (XZ 平面) ---
    \draw[->, thick, gray!80!black] (-2.2, 0) -- (2.3, 0) node[right] {$X$};
    \draw[->, thick, gray!80!black] (0, -1.6) -- (0, 2.6) node[above] {$Z$};
    \node[below left=2pt] at (O) {$O$};

    % --- 5. 中心・接点・頂点のプロットと寸法線 ---
    % 内接点 A(0, -1)
    \fill[red!85!black] (A) circle (1.3pt);
    \node[below right=1pt, font=\scriptsize, red!90!black] at (A) {$A(0, -1)$（内接点）};

    % 外側円の中心 M(0, 1/2)
    \fill[blue!90!black] (M) circle (1.2pt);
    \node[left=2pt, font=\scriptsize, blue!90!black] at (M) {$\frac{1}{2}$};

    % 外側円の半径 3/2 (矢印表示)
    \draw[<->,line width=1pt, teal!85!black] (M) -- (35:1.7)
    node[midway, above left=-1pt, font=\scriptsize] {$\frac{3}{2}$};

    % 内側円の半径 1 (矢印表示)
    \draw[<->, line width=1pt, red!80!black] (O) -- (145:1.0)
    node[midway, below left=-1pt, font=\scriptsize] {$1$};

    % 軸の目盛り
    \node[left=2pt, font=\scriptsize] at (0, 1.0) {$1$};
    \node[left=2pt, font=\scriptsize] at (0, 2.0) {$2$};
    \node[below=2pt, font=\scriptsize] at (1.0, 0) {$1$};
    \node[below=2pt, font=\scriptsize] at (-1.0, 0) {$-1$};
    \node[below=2pt, font=\scriptsize] at (1.414, 0) {$\sqrt{2}$}; % X 切片: X^2 + 1/4 = 9/4 => X = sqrt(2)
    \fill (1.414, 0) circle (0.6pt);

  \end{tikzpicture}
  \caption{領域 $V$ の $XZ$ 平面における切断面（点 $A$ で内接する 2 つの球の差）}
  \label{fig:xz_cross_section_inscribed_spheres}
\end{figure}

この体積$V$は$S$が半径$3/2$の球（中心$(0,0,1/2)$）に点$A$で内接し完全に含まれることから，
この二つの球の体積の差で
\begin{align}
  V
  =\dfrac{4}{3}\pi\left\{\left(\dfrac{3}{2}\right)^3-1\right\}=\dfrac{19}{6}\pi
  <\dfrac{19}{6}\times3.15=\dfrac{19.95}{2}<10
\end{align}
となる．

\bigskip

\end{document}
